\chapter{Registers and Counters}
\label{ch:regs}

From the previous chapter we know that one flip-flop stores one bit, and that
the security panel's alarm can be remembered and cleared with a key. But the
panel needs more: to store the lock code, to \emph{count how many seconds the
siren sounds}, and to make it ``beep'' once a second, even though the board's
clock generator runs at 27~MHz. The question of this chapter: what happens if
we connect many flip-flops, and how can we build storage, shifting and counting
from them?

\begin{buildbox}
In Logra, open Learn, the ``Counters and dividers'' chapter of the ``Circuits and
the board'' book, and go through the lesson ``A counter as one block''. Then the
lesson ``A divide-by-2 frequency divider''. In the \term{Waveforms} tab look at
all the counter bits at the same time: each next one changes half as often as the
previous one. Write down at which number the counter returns to zero and why.
Below we will see how this works.
\end{buildbox}

\newcommand{\dff}[3]{% name, x, y  -- D flip-flop 12x14 mm
  \draw[thick, fill=green!8] (#2,#3) rectangle ++(1.2,1.4);
  \node[font=\sffamily\scriptsize] at (#2+0.22,#3+1.05) {D};
  \node[font=\sffamily\scriptsize] at (#2+0.98,#3+1.05) {Q};
  \draw[thick] (#2,#3+0.2) -- (#2+0.18,#3+0.35) -- (#2,#3+0.5);
  \coordinate (#1-d) at (#2,#3+1.05); \coordinate (#1-q) at (#2+1.2,#3+1.05);
  \coordinate (#1-c) at (#2,#3+0.35);}

\section{The Storage Register}

A \term{register} is a group of D flip-flops with a common clock signal. At the
edge all the flip-flops simultaneously store their inputs, and the register
holds an $n$-bit number until the next edge. In schematics a register is
labeled with the letters RG.

\begin{figure}[H]
\centering
\begin{tikzpicture}
  \foreach \i in {0,...,3} {
    \pgfmathsetmacro{\x}{(3-\i)*2.3}
    \dff{r\i}{\x}{0}
    \draw[wire] (r\i-d) -- ++(-0.4,0) |- ($(r\i-d)+(-0.4,1.0)$) node[above] {$d_\i$};
    \draw[wire] (r\i-q) -- ++(0.35,0) -- ++(0,-1.5) node[below] {$q_\i$};
    \draw[wire, FpgaBlue] (r\i-c) -- ++(-0.3,0) -- ++(0,-0.75);
  }
  \draw[wire, FpgaBlue] (-1.4,-0.4) node[left, text=FpgaDark] {CLK} -- (6.6,-0.4);
  \foreach \i in {0,...,3} { \pgfmathsetmacro{\x}{(3-\i)*2.3-0.3} \fill[FpgaBlue] (\x,-0.4) circle (1.3pt); }
\end{tikzpicture}
\caption{A four-bit register: four D flip-flops with a common clock}
\label{fig:reg}
\end{figure}

\begin{figure}[H]
\centering
\begin{tikztimingtable}[timing/slope=0.08, xscale=1.5, yscale=1.4, timing/font=\sffamily\small, timing/rowdist=3.2ex]
  CLK      & 8{1L 1H} \\
  $D[3:0]$ & 2D{5} 1.5D{9} 1D{7} 1.5D{3} 4D{C} 6D{0} \\
  $Q[3:0]$ & 1D{?} 2D{5} 2D{9} 2D{3} 4D{C} 5D{0} \\
\end{tikztimingtable}
\caption{The register ``photographs'' the value of the bus $D$ at the moments of
the rising edges of CLK. The value 7 stayed on the input for less than one clock
cycle and fell between the edges---the register simply did not notice it}
\label{fig:reg_t}
\end{figure}

Registers are the ``variables'' of hardware. In a processor, registers hold
operands and results of computations; in any digital device they hold the
current state, counters, and data between processing stages. To make a register
change its value not on every edge but only on command, flip-flops with an
enable input CE are used (Fig.~\ref{fig:ce}).

\begin{diybox}[The Storage Register]
\textbf{Task.} \stars{1} The lesson ``A four-bit register''. \stars{2} The lesson
``A register with write enable'', then ``A register with reset'': load a number,
disable writing and make sure it does not change. \stars{3} Store the lock code
(four bits) in a register and show it on four LEDs.\\
\textbf{In Logra:} Learn, Designer, Waveforms.\\
\textbf{Expected result:} the output updates only on the edge, and only
if writing is enabled.\\
\textbf{How to check:} in Waveforms the value 7 on the input that arrived between
edges did not appear at the output (as in Fig.~\ref{fig:reg_t}).
\end{diybox}

\section{The Shift Register}

Connect the flip-flops in a chain: the output of each feeds the input of the
next. At every edge the contents \term{shift} by one bit, and a new bit from the
input $D_\text{in}$ is written into the first flip-flop.

\begin{figure}[H]
\centering
\begin{tikzpicture}
  \foreach \i in {0,...,3} {
    \pgfmathsetmacro{\x}{\i*2.3}
    \dff{s\i}{\x}{0}
    \draw[wire, FpgaBlue] (s\i-c) -- ++(-0.3,0) -- ++(0,-0.75);
  }
  \draw[wire] (-1.0,1.05) node[left] {$D_\text{in}$} -- (s0-d);
  \foreach \i [evaluate=\i as \j using int(\i+1)] in {0,1,2} {
    \draw[wire] (s\i-q) -- (s\j-d);
    \fill ($(s\i-q)+(0.3,0)$) circle (1.3pt);
    \draw[wire] ($(s\i-q)+(0.3,0)$) -- ++(0,1.0) node[above] {$q_\i$};
  }
  \draw[wire] (s3-q) -- ++(0.7,0) node[right] {$q_3 = D_\text{out}$};
  \draw[wire, FpgaBlue] (-1.0,-0.4) node[left, text=FpgaDark] {CLK} -- (6.6,-0.4);
\end{tikzpicture}
\caption{A four-bit shift register}
\label{fig:shift}
\end{figure}

\begin{figure}[H]
\centering
\begin{tikztimingtable}[timing/slope=0.08, xscale=1.5, yscale=1.4, timing/font=\sffamily\small, timing/rowdist=3.2ex]
  CLK            & 8{1L 1H} \\
  $D_\text{in}$  & 2H 2L 12L \\
  $q_0$          & 1L 2H 13L \\
  $q_1$          & 3L 2H 11L \\
  $q_2$          & 5L 2H 9L \\
  $q_3$          & 7L 2H 7L \\
\end{tikztimingtable}
\caption{A one applied to the input advances through the register by one bit
per clock cycle}
\end{figure}

\subsection{Why Shift}

\begin{description}
  \item[Serial $\to$ parallel code.] Bits arrive one at a time on a single wire,
    one after another (this is how data is sent over long cables and
    wireless channels: fewer wires). After
    $n$ clock cycles all $n$ bits are in the register at the same time and can be
    read as a number. \lookahead{this is how serial
    interfaces work, Chapter~\ref{ch:uart}.}
  \item[Parallel $\to$ serial code.] The reverse task when transmitting: the
    number is loaded into the register all at once and then ``pushed out'' one
    bit per clock cycle. For loading, a ``shift/load'' multiplexer is placed in
    front of each flip-flop.
  \item[Delaying a signal] by a given number of clock cycles.
  \item[Multiplying and dividing by 2]: shifting a binary number left by one bit
    multiplies it by 2, shifting right divides it.
\end{description}

\begin{figure}[H]
\centering
\begin{tikzpicture}[bit/.style={draw=FpgaBlue, fill=FpgaLight, minimum size=8mm, font=\ttfamily}]
  \node[font=\sffamily\small, anchor=east] at (-0.6,0) {transmitter};
  \foreach \b [count=\i from 0] in {1,0,1,1} \node[bit] at (\i*0.8,0) {\b};
  \draw[very thick, FpgaOrange, -{Stealth}] (3.0,0) -- node[above, font=\sffamily\small, text=FpgaDark] {1 wire: \code{1}, \code{1}, \code{0}, \code{1}\ldots} (7.4,0);
  \foreach \b [count=\i from 0] in {1,0,1,1} \node[bit] at (7.8+\i*0.8,0) {\b};
  \node[font=\sffamily\small, anchor=west] at (10.8,0) {receiver};
  \node[lbl] at (1.2,-0.8) {parallel $\to$ serial};
  \node[lbl] at (9.0,-0.8) {serial $\to$ parallel};
\end{tikzpicture}
\caption{Sending a number over one wire with two shift registers}
\end{figure}

\subsection{The Ring Register and the Pseudorandom Number Generator}

If the output of the last flip-flop is connected to the input of the first, you
get a \term{ring register}: a one written into it will run around the circle
forever, like the ``running light'' of a garland. \lookahead{a ready-made
``running lights'' example is in Chapter~\ref{ch:examples}.}

And if you feed the input with the XOR of several bits, you get a
\term{linear feedback shift register} (LFSR). With a proper choice of
taps it goes through all $2^n-1$ nonzero states in a ``shuffled''
order. Such registers are used as simple pseudorandom number generators, in
error-correcting coding and for encryption.

\begin{figure}[H]
\centering
\begin{minipage}{0.55\textwidth}
\centering
\begin{tikzpicture}
  \foreach \i in {0,...,2} {
    \pgfmathsetmacro{\x}{\i*2.0}
    \dff{l\i}{\x}{0}
  }
  \draw[wire] (l0-q) -- (l1-d); \draw[wire] (l1-q) -- (l2-d);
  \node[gXOR, minimum width=9mm, minimum height=7mm, rotate=180] (x) at (2.6,2.3) {};
  \draw[wire] (l2-q) -- ++(0.4,0) |- (x.input 1);
  \fill ($(l1-q)+(0.3,0)$) circle (1.3pt);
  \draw[wire] ($(l1-q)+(0.3,0)$) |- (x.input 2);
  \draw[wire] (x.output) -| ($(l0-d)+(-0.4,0)$) -- (l0-d);
  \foreach \i in {0,...,2} \node[lbl, below=1mm] at ($(l\i-c)+(0.6,-0.4)$) {$q_\i$};
\end{tikzpicture}
\end{minipage}\hfill
\begin{minipage}{0.4\textwidth}
\centering
\small
\begin{tabular}{c|ccc}
\toprule
cycle & $q_0$ & $q_1$ & $q_2$ \\
\midrule
0 & 1 & 0 & 0 \\
1 & 0 & 1 & 0 \\
2 & 1 & 0 & 1 \\
3 & 1 & 1 & 0 \\
4 & 1 & 1 & 1 \\
5 & 0 & 1 & 1 \\
6 & 0 & 0 & 1 \\
7 & 1 & 0 & 0 \\
\bottomrule
\end{tabular}
\end{minipage}
\caption{A three-bit LFSR ($q_0 \leftarrow q_1\oplus q_2$) goes through all seven
nonzero states and repeats}
\label{fig:lfsr}
\end{figure}

The ``all zeros'' state is forbidden for such a register: from it
$q_1\oplus q_2=0$, and the register stays zero forever.

\begin{diybox}[Shift, Ring and LFSR]
\textbf{Task.} \stars{1} The lesson ``A shift register as a block''. \stars{2} The lessons
``Serial input'' and ``Assembling a word from a stream'': check in which order
(most or least significant bit first) the bits enter the word. \stars{3} The lessons
``Ring counter'', ``Johnson counter'' and ``Pseudorandom numbers: LFSR'': make sure that
a four-bit LFSR goes through 15 states.\\
\textbf{In Logra:} Learn, Designer, Waveforms; for enthusiasts, the
example \code{lfsr6}.\\
\textbf{Expected result:} the one advances by one bit per clock cycle;
the ring returns it to the starting position after $n$ cycles.\\
\textbf{How to check:} Waveforms: write out the states and compare with the table
in Fig.~\ref{fig:lfsr}.
\end{diybox}

\section{Counters}

A \term{counter} is a register whose value increases (or decreases) by one on
every clock cycle. Counters measure time, divide
frequency and step through memory addresses.

\subsection{The Asynchronous Counter}

Look at how the bits change when counting $0, 1, 2, \ldots, 7$: the lowest bit
toggles at every step, the next one half as often, the next one half as often
again. Each bit toggles when the previous one goes from 1
to 0. So T flip-flops can be chained: the clock input of each
next flip-flop is connected to the \emph{inverted} output of the previous one.

\begin{figure}[H]
\centering
\begin{tikzpicture}
  \foreach \i in {0,...,2} {
    \pgfmathsetmacro{\x}{\i*2.8}
    \draw[thick, fill=green!8] (\x,0) rectangle ++(1.2,1.4);
    \node[font=\sffamily\scriptsize] at (\x+0.22,1.05) {T};
    \node[font=\sffamily\scriptsize] at (\x+0.98,1.05) {Q};
    \node[font=\sffamily\scriptsize] at (\x+0.98,0.35) {$\overline{\text{Q}}$};
    \draw[thick] (\x,0.2) -- (\x+0.18,0.35) -- (\x,0.5);
    \draw[wire] (\x,1.05) -- ++(-0.3,0) node[left, font=\scriptsize] {1};
    \draw[wire] (\x+1.2,1.05) -- ++(0.4,0) -- ++(0,0.8) node[above] {$q_\i$};
  }
  \draw[wire, FpgaBlue] (-1.0,0.35) node[left, text=FpgaDark] {CLK} -- (0,0.35);
  \draw[wire, FpgaRed] (1.2,0.35) -- (2.8,0.35);
  \draw[wire, FpgaRed] (4.0,0.35) -- (5.6,0.35);
\end{tikzpicture}
\caption{An asynchronous (ripple) three-bit counter made of T flip-flops}
\label{fig:ripple}
\end{figure}

\begin{figure}[H]
\centering
\begin{tikztimingtable}[timing/slope=0.05, xscale=1.2, yscale=1.4, timing/font=\sffamily\small, timing/rowdist=3.2ex]
  CLK     & 16{1L 1H} \\
  $q_0$   & 1L 7{2H 2L} 2H 1L \\
  $q_1$   & 3L 3{4H 4L} 4H 1L \\
  $q_2$   & 7L 8H 8L 8H 1L \\
  value   & 1D{0} 2D{1} 2D{2} 2D{3} 2D{4} 2D{5} 2D{6} 2D{7} 2D{0} 2D{1} 2D{2} 2D{3} 2D{4} 2D{5} 2D{6} 2D{7} 1D{0} \\
\end{tikztimingtable}
\caption{Each next bit changes half as often: the counter steps through the numbers
$0\ldots7$, and the bit $q_k$ has the frequency $f_\text{CLK}/2^{k+1}$}
\label{fig:count_t}
\end{figure}

The drawback of the asynchronous counter is that the flip-flops switch not
simultaneously but in a ``wave'', each with a delay relative to the previous one.
At the transition $011 \to 100$ false values appear at the output for a short time
($011\to010\to000\to100$). In a long counter this ``wave'' takes a long time.
So an asynchronous counter is good only where its outputs are not needed by the
rest of the circuit at the same time (for example, as a frequency divider), and
never for comparing the value with a constant.

\subsection{The Synchronous Counter}

In a synchronous counter all the flip-flops are clocked by \emph{one} signal, and
the new value is computed by a combinational circuit---an adder that adds one.
All the bits change simultaneously on the edge. This is how almost all counters in
modern digital devices are built.

\begin{figure}[H]
\centering
\begin{tikzpicture}
  \node[draw=FpgaGreen, fill=green!8, thick, minimum width=27mm, minimum height=18mm, font=\sffamily\small, align=center] (reg) at (4.4,0) {register\\\scriptsize $n$ bits};
  \node[font=\sffamily\scriptsize, anchor=west] at ([xshift=0.5mm]reg.west) {D};
  \node[font=\sffamily\scriptsize, anchor=east] at ([xshift=-0.5mm]reg.east) {Q};
  \draw[thick] ([xshift=-1.2mm]reg.south) -- ([yshift=1.8mm]reg.south) -- ([xshift=1.2mm]reg.south);
  \draw[wire, FpgaBlue] (reg.south) -- ++(0,-0.8) node[below, text=FpgaDark] {CLK};
  \node[draw=FpgaBlue, fill=FpgaLight, thick, circle, minimum size=12mm, font=\sffamily\bfseries] (add) at (1.4,0) {$+1$};
  \draw[bus, -{Stealth}] (add.east) -- (reg.west);
  \draw[bus, -{Stealth}] (reg.east) -- ++(1.8,0) node[right, text=FpgaDark] {$Q$};
  \draw[bus, -{Stealth}] ($(reg.east)+(0.9,0)$) -- ++(0,1.6) -| (add.north);
  \fill[FpgaBlue] ($(reg.east)+(0.9,0)$) circle (1.8pt);
  \node[lbl] at (4.4,1.9) {feedback};
\end{tikzpicture}
\caption{A synchronous counter: a register plus an adder. Every clock cycle the register
stores $Q+1$}
\label{fig:sync_counter}
\end{figure}

\begin{diybox}[The Counter]
\textbf{Task.} \stars{1} The lesson ``A counter as one block'': what happens after 15?
\stars{2} The lessons ``A counter with pause'' and ``A counter from a register and
\,`+1'\,'': compare the synchronous counter built by hand with the ready-made
block. \stars{3} Build a counter that counts only while a button is pressed.\\
\textbf{In Logra:} Learn, Designer, Waveforms, RTL; the example
\code{counter}.\\
\textbf{Expected result:} after $1111$ the counter goes to $0000$.\\
\textbf{How to check:} RTL will show a register and a ``+1'' adder, as in
Fig.~\ref{fig:sync_counter}.
\end{diybox}

\section{The Frequency Divider and the Modulo Counter}

As Fig.~\ref{fig:count_t} shows, each bit of a counter is a frequency divider
by a power of two. But how do we divide by a number that is not a power of two,
for example by 10? We need a \term{modulo-$N$ counter}: it counts
$0, 1, \ldots, N-1$ and then resets to 0. For this a comparator is added to the
circuit, which compares the value with $N-1$, and a multiplexer, which feeds 0 to
the register instead of $Q+1$. The comparator output $Q=N-1$ is a pulse
one clock cycle long, and it appears exactly once every $N$ cycles: this is how
the ``one second has passed'' signal is obtained from fast clock cycles.

\begin{figure}[H]
\centering
\begin{tikzpicture}
  \node[draw=FpgaGreen, fill=green!8, thick, minimum width=18mm, minimum height=16mm, font=\sffamily\small] (reg) at (5.6,0) {register};
  \draw[wire, FpgaBlue] (reg.south) -- ++(0,-0.6) node[right, text=FpgaDark, font=\sffamily\small] {CLK};
  \draw[thick, fill=orange!15] (3.0,0.9) -- (3.8,0.5) -- (3.8,-0.5) -- (3.0,-0.9) -- cycle;
  \node[font=\sffamily\scriptsize] at (3.4,0) {MUX};
  \node[font=\tiny, anchor=west] at (3.02,0.5) {0}; \node[font=\tiny, anchor=west] at (3.02,-0.5) {1};
  \node[draw=FpgaBlue, fill=FpgaLight, thick, circle, minimum size=10mm, font=\sffamily\bfseries\small] (add) at (1.2,0.5) {$+1$};
  \draw[bus, -{Stealth}] (add.east) -- (3.0,0.5);
  \draw[bus, -{Stealth}] (1.8,-0.5) node[left, text=FpgaDark] {0} -- (3.0,-0.5);
  \draw[bus, -{Stealth}] (3.8,0) -- (reg.west);
  \draw[bus, -{Stealth}] (reg.east) -- ++(1.6,0) node[right, text=FpgaDark] {$Q$};
  \draw[bus] ($(reg.east)+(0.8,0)$) -- ++(0,1.6) -| (add.north);
  \fill[FpgaBlue] ($(reg.east)+(0.8,0)$) circle (1.8pt);
  \node[block, minimum width=20mm] (cmp) at (5.6,-2.9) {$Q = N-1$?};
  \draw[bus, -{Stealth}] ($(reg.east)+(0.8,0)$) |- (cmp.east);
  \draw[arr] (cmp.west) -| node[pos=0.3, below, lbl] {reset} (3.4,-0.72);
\end{tikzpicture}
\caption{A modulo-$N$ counter}
\end{figure}

\begin{figure}[H]
\centering
\begin{tikzpicture}
\begin{axis}[
  width=0.82\textwidth, height=5.6cm,
  xlabel={Counter width $n$}, ylabel={Frequency, Hz},
  ymode=log, ymin=0.1, ymax=3e7, xmin=0, xmax=26, xtick={1,4,8,12,16,20,25},
  label style={font=\sffamily\small}, tick label style={font=\sffamily\small},
  axis lines*=left, grid=major, grid style={FpgaGray!20},
]
  \addplot[very thick, FpgaBlue, mark=*, mark size=1.5pt, samples at={1,2,...,25}] {27e6/2^x};
  \addplot[FpgaRed, thick, dashed, domain=0:26, samples=2] {1};
  \node[font=\sffamily\small, text=FpgaRed, anchor=south west] at (axis cs:1,1.3) {1 Hz---blinking visible to the eye};
  \node[font=\sffamily\small, text=FpgaBlue, anchor=south west] at (axis cs:25,0.8) {$n=25$};
\end{axis}
\end{tikzpicture}
\caption{The frequency at the output of the most significant bit of an $n$-bit counter with
a 27\,MHz clock: $f = 27\,\text{MHz}/2^n$. To get blinking at about 1~Hz, you need a
counter of about 25 bits}
\label{fig:div_plot}
\end{figure}

But the output of the most significant bit of a 25-bit counter gives not 1~Hz but
$27\cdot10^6/2^{25}\approx0.8$~Hz. To get exactly 1~Hz you need a counter
modulo $N=27\,000\,000$: it counts from 0 to $26\,999\,999$ (this is also 25
bits) and once a second produces a pulse. If this pulse toggles a
T flip-flop, you get a 0.5~Hz square wave (one period in two seconds). For a
square wave of exactly 1~Hz, take the modulus $13\,500\,000$: the pulse arrives every
half second and toggles the flip-flop.

\begin{diybox}[Modulus $N$ and the Frequency Divider; Project Stage]
\textbf{Task.} \stars{1} The lessons ``Counting from 0 to 5'' (modulo 6) and ``A divide-by-4
from a counter''. \stars{2} The lessons ``A counter with reset'' and ``A counter with load'',
then ``A counter to a seven-segment display'' (chapter ``Projects'' of the ``Circuits
and the board'' book): make a modulo-10 counter shown on the display.
\stars{3} \textbf{Project stage ``Security Panel'': the alarm-time counter and the
``beep''.} Build a counter that counts seconds while the alarm from stage~7
is on and stops at 9. Get the seconds with a divider: so that
the check in Waveforms does not take 27 million clock cycles, first set a
small modulus (for example, $N=4$), and leave the real $N=27\,000\,000$ for the board.
The buzzer ``beeps'' from the seconds bit.\\
\textbf{In Logra:} Learn (chapters ``Counters and dividers'', ``Projects''), Designer, Waveforms,
RTL; the examples \code{blink}, \code{countdown}, \code{stopwatch}.\\
\textbf{Expected result:} with $N=4$ the ``second'' pulse arrives once every 4
clock cycles; the seconds counter grows $0,1,\ldots,9$ and freezes; on a reset by the key it
is cleared.\\
\textbf{Acceptance criterion:} in Waveforms measure the pulse period
(it equals $N$ clock cycles), the alarm output does not grow beyond the 9th second,
and on reset all counters are zero on the very next clock cycle.
\end{diybox}

\begin{ideabox}
A \textbf{register} is a group of flip-flops with a common clock; it stores a number.
A \textbf{shift register} moves bits along a chain and converts
serial code to parallel and back. A \textbf{synchronous counter}
is a register with an adder in the feedback path; its bits divide the clock frequency
by powers of two.
\end{ideabox}

\section{Common Mistakes}

\begin{warnbox}[Common Mistakes]
\begin{enumerate}
  \item \textbf{An off-by-one error in the modulus.} A modulo-$N$ counter is reset
        at the value $N-1$, not $N$. If you compare with $N$, there will be
        $N+1$ states, and the extra value $N$ will appear at the output for an
        instant.
  \item \textbf{``A bit flies through the whole register in one clock cycle.''} All the flip-flops
        shift simultaneously and take the \emph{old} values of their neighbors, so
        a bit advances exactly one position per edge.
  \item \textbf{A mixed-up bit order.} When receiving a stream, the most significant bit
        may come first or last; if the word is assembled by shifting in the
        other direction, you get a number with the reverse bit order
        ($0011$ becomes $1100$, and $0001$ becomes $1000$).
  \item \textbf{Confusing ``blinking'' and ``pulse''.} The most significant bit of a counter
        gives a square wave, whereas the comparison $Q=N-1$ gives one short pulse per cycle.
        To toggle another node you need exactly a pulse, not a square wave.
  \item \textbf{The zero state of an LFSR.} A register with an XOR that is filled
        with zeros will never leave them; the initial value must be
        nonzero.
  \item \textbf{Using the outputs of an asynchronous counter as data.}
        Because of the switching ``wave'', false numbers flicker on the outputs, and logic
        that reads them without synchronization may fire an extra time.
\end{enumerate}
\end{warnbox}

\section{Questions and Problems}
\label{sec:tasks-regs}

Answers are collected in Appendix~\ref{app:answers}.

\begin{taskbox}[Chapter Problems]
\begin{enumerate}
  \item \stars{1}~A four-bit ring register holds $0001$. What values will it take over the next four clock cycles?
  \item \stars{1}~A four-bit binary counter is in the state $1111$. What will its state be after the next clock cycle? What is the counting modulus?
  \item \stars{2}~How many bits must a counter have to divide 27~MHz down to 1~Hz, if it counts up to 27\,000\,000?
  \item \stars{1}~What value does a modulo-10 counter count up to and how is it reset?
  \item \stars{2}~How many different nonzero states does a 4-bit LFSR with the feedback $x^4 + x^3 + 1$ pass through before repeating?
  \item \stars{1}~What is the main advantage of a synchronous counter over an asynchronous one?
  \item \stars{2}~\textbf{Find the error.} A student built a modulo-10 counter and reset it at the value $1010$. What sequence does the output see? What should be fixed?
  \item \stars{2}~\textbf{Find the error.} The three-bit LFSR of Fig.~\ref{fig:lfsr} was powered on without setting an initial value and ended up in the state $000$. What will happen next and how can it be avoided?
  \item \stars{3}~\textbf{Project.} A modulo-$N=6$ counter (or any other $N$ of your choice) showing the number on a seven-segment display. Acceptance criterion: Waveforms shows the cycle $0,1,\ldots,N-1,0$, the value $N$ does not appear at any moment, and the ``carry'' output is one pulse per cycle at the number $N-1$.
  \item \stars{3}~\textbf{Project (stage 8).} A divider from 27~MHz to 1~Hz: first with a small modulus in Waveforms, then with the real one. Acceptance criterion: for the modulus $N$ the period of the output pulse equals $N$ clock cycles; for the board the modulus is $N=27\,000\,000$ and the counter has 25 bits.
  \item \stars{3}~\textbf{Project (stage 8).} An alarm-time counter: it counts seconds only while the alarm is remembered (stage~7), stops at 9 and is cleared by the key. Acceptance criterion: four checks in Waveforms---growth, stop at 9, stop when the alarm is cleared, clearing by the key.
\end{enumerate}
\end{taskbox}

\section{What's Next}

Now the panel has memory, numbers and time: the lock code in a register, the alarm
seconds in a counter and the ``beep'' from a divider. But all of this is still
assembled by hand, and the circuit has no ``modes'': the key is turned on, the signal
goes, the alarm turns on, the code is entered. How can we describe such a sequence
of states rigorously, so that nothing is lost? For that we need synchronous circuits
and finite-state machines---the subject of the next chapter.
